\IfFileExists{stacks-project.cls}{%
\documentclass{stacks-project}
}{%
\documentclass{amsart}
}

% For dealing with references we use the comment environment
\usepackage{verbatim}
\newenvironment{reference}{\comment}{\endcomment}
%\newenvironment{reference}{}{}
\newenvironment{slogan}{\comment}{\endcomment}
\newenvironment{history}{\comment}{\endcomment}

% For commutative diagrams we use Xy-pic
\usepackage[all]{xy}

% We use 2cell for 2-commutative diagrams.
\xyoption{2cell}
\UseAllTwocells

% We use multicol for the list of chapters between chapters
\usepackage{multicol}

% This is generally recommended for better output
\usepackage{lmodern}
\usepackage[T1]{fontenc}

% For cross-file-references
\usepackage{xr-hyper}

% Package for hypertext links:
\usepackage{hyperref}

% For any local file, say "hello.tex" you want to link to please
% use \externaldocument[hello-]{hello}
\externaldocument[introduction-]{introduction}
\externaldocument[conventions-]{conventions}
\externaldocument[sets-]{sets}
\externaldocument[categories-]{categories}
\externaldocument[topology-]{topology}
\externaldocument[sheaves-]{sheaves}
\externaldocument[sites-]{sites}
\externaldocument[stacks-]{stacks}
\externaldocument[fields-]{fields}
\externaldocument[algebra-]{algebra}
\externaldocument[brauer-]{brauer}
\externaldocument[homology-]{homology}
\externaldocument[derived-]{derived}
\externaldocument[simplicial-]{simplicial}
\externaldocument[more-algebra-]{more-algebra}
\externaldocument[smoothing-]{smoothing}
\externaldocument[modules-]{modules}
\externaldocument[sites-modules-]{sites-modules}
\externaldocument[injectives-]{injectives}
\externaldocument[cohomology-]{cohomology}
\externaldocument[sites-cohomology-]{sites-cohomology}
\externaldocument[dga-]{dga}
\externaldocument[dpa-]{dpa}
\externaldocument[sdga-]{sdga}
\externaldocument[hypercovering-]{hypercovering}
\externaldocument[schemes-]{schemes}
\externaldocument[constructions-]{constructions}
\externaldocument[properties-]{properties}
\externaldocument[morphisms-]{morphisms}
\externaldocument[coherent-]{coherent}
\externaldocument[divisors-]{divisors}
\externaldocument[limits-]{limits}
\externaldocument[varieties-]{varieties}
\externaldocument[topologies-]{topologies}
\externaldocument[descent-]{descent}
\externaldocument[perfect-]{perfect}
\externaldocument[more-morphisms-]{more-morphisms}
\externaldocument[flat-]{flat}
\externaldocument[groupoids-]{groupoids}
\externaldocument[more-groupoids-]{more-groupoids}
\externaldocument[etale-]{etale}
\externaldocument[chow-]{chow}
\externaldocument[intersection-]{intersection}
\externaldocument[pic-]{pic}
\externaldocument[weil-]{weil}
\externaldocument[adequate-]{adequate}
\externaldocument[dualizing-]{dualizing}
\externaldocument[duality-]{duality}
\externaldocument[discriminant-]{discriminant}
\externaldocument[derham-]{derham}
\externaldocument[local-cohomology-]{local-cohomology}
\externaldocument[algebraization-]{algebraization}
\externaldocument[curves-]{curves}
\externaldocument[resolve-]{resolve}
\externaldocument[models-]{models}
\externaldocument[functors-]{functors}
\externaldocument[equiv-]{equiv}
\externaldocument[pione-]{pione}
\externaldocument[etale-cohomology-]{etale-cohomology}
\externaldocument[proetale-]{proetale}
\externaldocument[relative-cycles-]{relative-cycles}
\externaldocument[more-etale-]{more-etale}
\externaldocument[trace-]{trace}
\externaldocument[crystalline-]{crystalline}
\externaldocument[spaces-]{spaces}
\externaldocument[spaces-properties-]{spaces-properties}
\externaldocument[spaces-morphisms-]{spaces-morphisms}
\externaldocument[decent-spaces-]{decent-spaces}
\externaldocument[spaces-cohomology-]{spaces-cohomology}
\externaldocument[spaces-limits-]{spaces-limits}
\externaldocument[spaces-divisors-]{spaces-divisors}
\externaldocument[spaces-over-fields-]{spaces-over-fields}
\externaldocument[spaces-topologies-]{spaces-topologies}
\externaldocument[spaces-descent-]{spaces-descent}
\externaldocument[spaces-perfect-]{spaces-perfect}
\externaldocument[spaces-more-morphisms-]{spaces-more-morphisms}
\externaldocument[spaces-flat-]{spaces-flat}
\externaldocument[spaces-groupoids-]{spaces-groupoids}
\externaldocument[spaces-more-groupoids-]{spaces-more-groupoids}
\externaldocument[bootstrap-]{bootstrap}
\externaldocument[spaces-pushouts-]{spaces-pushouts}
\externaldocument[spaces-chow-]{spaces-chow}
\externaldocument[groupoids-quotients-]{groupoids-quotients}
\externaldocument[spaces-more-cohomology-]{spaces-more-cohomology}
\externaldocument[spaces-simplicial-]{spaces-simplicial}
\externaldocument[spaces-duality-]{spaces-duality}
\externaldocument[formal-spaces-]{formal-spaces}
\externaldocument[restricted-]{restricted}
\externaldocument[spaces-resolve-]{spaces-resolve}
\externaldocument[formal-defos-]{formal-defos}
\externaldocument[defos-]{defos}
\externaldocument[cotangent-]{cotangent}
\externaldocument[examples-defos-]{examples-defos}
\externaldocument[algebraic-]{algebraic}
\externaldocument[examples-stacks-]{examples-stacks}
\externaldocument[stacks-sheaves-]{stacks-sheaves}
\externaldocument[criteria-]{criteria}
\externaldocument[artin-]{artin}
\externaldocument[quot-]{quot}
\externaldocument[stacks-properties-]{stacks-properties}
\externaldocument[stacks-morphisms-]{stacks-morphisms}
\externaldocument[stacks-limits-]{stacks-limits}
\externaldocument[stacks-cohomology-]{stacks-cohomology}
\externaldocument[stacks-perfect-]{stacks-perfect}
\externaldocument[stacks-introduction-]{stacks-introduction}
\externaldocument[stacks-more-morphisms-]{stacks-more-morphisms}
\externaldocument[stacks-geometry-]{stacks-geometry}
\externaldocument[moduli-]{moduli}
\externaldocument[moduli-curves-]{moduli-curves}
\externaldocument[examples-]{examples}
\externaldocument[exercises-]{exercises}
\externaldocument[guide-]{guide}
\externaldocument[desirables-]{desirables}
\externaldocument[coding-]{coding}
\externaldocument[obsolete-]{obsolete}
\externaldocument[fdl-]{fdl}
\externaldocument[index-]{index}

% Theorem environments.
%
\theoremstyle{plain}
\newtheorem{theorem}[subsection]{Theorem}
\newtheorem{proposition}[subsection]{Proposition}
\newtheorem{lemma}[subsection]{Lemma}

\theoremstyle{definition}
\newtheorem{definition}[subsection]{Definition}
\newtheorem{example}[subsection]{Example}
\newtheorem{exercise}[subsection]{Exercise}
\newtheorem{situation}[subsection]{Situation}

\theoremstyle{remark}
\newtheorem{remark}[subsection]{Remark}
\newtheorem{remarks}[subsection]{Remarks}

\numberwithin{equation}{subsection}

% Macros
%
\def\lim{\mathop{\mathrm{lim}}\nolimits}
\def\colim{\mathop{\mathrm{colim}}\nolimits}
\def\Spec{\mathop{\mathrm{Spec}}}
\def\Hom{\mathop{\mathrm{Hom}}\nolimits}
\def\Ext{\mathop{\mathrm{Ext}}\nolimits}
\def\SheafHom{\mathop{\mathcal{H}\!\mathit{om}}\nolimits}
\def\SheafExt{\mathop{\mathcal{E}\!\mathit{xt}}\nolimits}
\def\Sch{\mathit{Sch}}
\def\Mor{\mathop{\mathrm{Mor}}\nolimits}
\def\Ob{\mathop{\mathrm{Ob}}\nolimits}
\def\Sh{\mathop{\mathit{Sh}}\nolimits}
\def\NL{\mathop{N\!L}\nolimits}
\def\CH{\mathop{\mathrm{CH}}\nolimits}
\def\proetale{{pro\text{-}\acute{e}tale}}
\def\etale{{\acute{e}tale}}
\def\QCoh{\mathit{QCoh}}
\def\Ker{\mathop{\mathrm{Ker}}}
\def\Im{\mathop{\mathrm{Im}}}
\def\Coker{\mathop{\mathrm{Coker}}}
\def\Coim{\mathop{\mathrm{Coim}}}

% Boxtimes
%
\DeclareMathSymbol{\boxtimes}{\mathbin}{AMSa}{"02}

%
% Macros for moduli stacks/spaces
%
\def\QCohstack{\mathcal{QC}\!\mathit{oh}}
\def\Cohstack{\mathcal{C}\!\mathit{oh}}
\def\Spacesstack{\mathcal{S}\!\mathit{paces}}
\def\Quotfunctor{\mathrm{Quot}}
\def\Hilbfunctor{\mathrm{Hilb}}
\def\Curvesstack{\mathcal{C}\!\mathit{urves}}
\def\Polarizedstack{\mathcal{P}\!\mathit{olarized}}
\def\Complexesstack{\mathcal{C}\!\mathit{omplexes}}
% \Pic is the operator that assigns to X its picard group, usage \Pic(X)
% \Picardstack_{X/B} denotes the Picard stack of X over B
% \Picardfunctor_{X/B} denotes the Picard functor of X over B
\def\Pic{\mathop{\mathrm{Pic}}\nolimits}
\def\Picardstack{\mathcal{P}\!\mathit{ic}}
\def\Picardfunctor{\mathrm{Pic}}
\def\Deformationcategory{\mathcal{D}\!\mathit{ef}}
\usepackage{cleveref}
\crefname{originalreference}{original reference}{original references}
\Crefname{originalreference}{Original reference}{Original references}
\makeatletter
\newlabel{categories-lemma-representable-over-representable}{{03KC}{source}{}{sourcepointer-1}{}}
\newlabel{categories-lemma-representable-over-representable@cref}{{[originalreference][0][]03KC}{[1][0][]source}{}{}{}}
\makeatother
\crefname{originallemma}{original lemma}{original lemmas}
\Crefname{originallemma}{Original lemma}{Original lemmas}
\makeatletter
\newlabel{lemma-U-universally-bounded}{{03JT}{source}{}{sourcepointer-2}{}}
\newlabel{lemma-U-universally-bounded@cref}{{[originallemma][0][]03JT}{[1][0][]source}{}{}{}}
\makeatother
\crefname{originallemma}{original lemma}{original lemmas}
\Crefname{originallemma}{Original lemma}{Original lemmas}
\makeatletter
\newlabel{lemma-UR-finite-above-x}{{03JV}{source}{}{sourcepointer-3}{}}
\newlabel{lemma-UR-finite-above-x@cref}{{[originallemma][0][]03JV}{[1][0][]source}{}{}{}}
\makeatother
\crefname{originallemma}{original lemma}{original lemmas}
\Crefname{originallemma}{Original lemma}{Original lemmas}
\makeatletter
\newlabel{lemma-base-change-universally-bounded}{{03JN}{source}{}{sourcepointer-4}{}}
\newlabel{lemma-base-change-universally-bounded@cref}{{[originallemma][0][]03JN}{[1][0][]source}{}{}{}}
\makeatother
\crefname{originallemma}{original lemma}{original lemmas}
\Crefname{originallemma}{Original lemma}{Original lemmas}
\makeatletter
\newlabel{lemma-bounded-fibres}{{03JX}{source}{}{sourcepointer-5}{}}
\newlabel{lemma-bounded-fibres@cref}{{[originallemma][0][]03JX}{[1][0][]source}{}{}{}}
\makeatother
\crefname{originallemma}{original lemma}{original lemmas}
\Crefname{originallemma}{Original lemma}{Original lemmas}
\makeatletter
\newlabel{lemma-characterize-very-reasonable}{{03IH}{source}{}{sourcepointer-6}{}}
\newlabel{lemma-characterize-very-reasonable@cref}{{[originallemma][0][]03IH}{[1][0][]source}{}{}{}}
\makeatother
\crefname{originalreference}{original reference}{original references}
\Crefname{originalreference}{Original reference}{Original references}
\makeatletter
\newlabel{spaces-lemma-base-change-representable-transformations-property}{{02WL}{source}{}{sourcepointer-7}{}}
\newlabel{spaces-lemma-base-change-representable-transformations-property@cref}{{[originalreference][0][]02WL}{[1][0][]source}{}{}{}}
\makeatother
\title{The point conditions beta, gamma, delta, epsilon and theta pull back along representable morphisms}
\author{The Stacks Project Authors}
\date{stacks-project:decent-spaces.tex:local-snapshot-20261012; GNU FDL 1.2 or later}
\begin{document}
\maketitle
\noindent\textit{Editorial note.} The theta case is a one-line citation of a Categories lemma (cross-chapter, named in the prerequisites); the base-change boundedness input is not carded and is declared external.\par\medskip
\noindent\textit{Allowed standard prerequisites (editorial).}\begin{itemize}
\item UR-finiteness characterization. lemma-UR-finite-above-x (carded L26700). Condition (2) expands (gamma) into the covering-family form.
\item U-universally-bounded characterization. lemma-U-universally-bounded (carded L26701). Closes the delta case.
\item Base change of universal boundedness. lemma-base-change-universally-bounded (same chapter, NOT carded, declared external). Gives bounded fibres of V x\_Y X -> X in the delta case.
\item Very-reasonable characterization. lemma-characterize-very-reasonable (carded L26702). Expands (epsilon) into the covering-family form.
\item Base-change transfer for representable transformations. Spaces, lemma base-change-representable-transformations-property (cross-chapter, named). Used throughout the epsilon case.
\item Representability over representable. Categories, lemma representable-over-representable (cross-chapter, named). The theta case.
\end{itemize}\medskip
\ifdefined\hypertarget\hypertarget{sourcepointer-1}{}\fi
\noindent\textit{Source reference.} Original reference 03KC: focused-seeds/geometry/S9-decent-representable-properties/ref-tags.txt, source lines 5--5. Navigation pointer only: the counted proof cites this statement by \textbackslash{}ref as a standing input; the deduction the counted proof has to make is carried out inside the counted span, and no part of that deduction is delegated to this pointer.\par\medskip
\ifdefined\hypertarget\hypertarget{sourcepointer-2}{}\fi
\noindent\textit{Source reference.} Original lemma 03JT: sources/stacks-decent-spaces/pinned-original-source.tex, source lines 847--847. Navigation pointer only: the counted proof cites this statement by \textbackslash{}ref as a standing input; the deduction the counted proof has to make is carried out inside the counted span, and no part of that deduction is delegated to this pointer.\par\medskip
\ifdefined\hypertarget\hypertarget{sourcepointer-3}{}\fi
\noindent\textit{Source reference.} Original lemma 03JV: sources/stacks-decent-spaces/pinned-original-source.tex, source lines 732--732. Navigation pointer only: the counted proof cites this statement by \textbackslash{}ref as a standing input; the deduction the counted proof has to make is carried out inside the counted span, and no part of that deduction is delegated to this pointer.\par\medskip
\ifdefined\hypertarget\hypertarget{sourcepointer-4}{}\fi
\noindent\textit{Source reference.} Original lemma 03JN: sources/stacks-decent-spaces/pinned-original-source.tex, source lines 327--327. Navigation pointer only: the counted proof cites this statement by \textbackslash{}ref as a standing input; the deduction the counted proof has to make is carried out inside the counted span, and no part of that deduction is delegated to this pointer.\par\medskip
\ifdefined\hypertarget\hypertarget{sourcepointer-5}{}\fi
\noindent\textit{Source reference.} Original lemma 03JX: sources/stacks-decent-spaces/pinned-original-source.tex, source lines 952--952. Navigation pointer only: the counted proof cites this statement by \textbackslash{}ref as a standing input; the deduction the counted proof has to make is carried out inside the counted span, and no part of that deduction is delegated to this pointer.\par\medskip
\ifdefined\hypertarget\hypertarget{sourcepointer-6}{}\fi
\noindent\textit{Source reference.} Original lemma 03IH: sources/stacks-decent-spaces/pinned-original-source.tex, source lines 896--896. Navigation pointer only: the counted proof cites this statement by \textbackslash{}ref as a standing input; the deduction the counted proof has to make is carried out inside the counted span, and no part of that deduction is delegated to this pointer.\par\medskip
\ifdefined\hypertarget\hypertarget{sourcepointer-7}{}\fi
\noindent\textit{Source reference.} Original reference 02WL: focused-seeds/geometry/S9-decent-representable-properties/ref-tags.txt, source lines 6--6. Navigation pointer only: the counted proof cites this statement by \textbackslash{}ref as a standing input; the deduction the counted proof has to make is carried out inside the counted span, and no part of that deduction is delegated to this pointer.\par\medskip
\begin{lemma}
\label{lemma-representable-properties}
Let $S$ be a scheme. Let $\mathcal{P}$ be one of the properties
$(\beta)$, $(\gamma)$, $(\delta)$, $(\epsilon)$, or
$(\theta)$ of algebraic spaces listed in
Lemma \ref{lemma-bounded-fibres}.
Let $X$, $Y$ be algebraic spaces over $S$.
Let $X \to Y$ be a representable morphism.
If $Y$ has property $\mathcal{P}$, so does $X$.
\end{lemma}

\begin{proof}
Assume $f : X \to Y$ is a representable morphism of algebraic spaces,
and assume that $Y$ has $\mathcal{P}$. Let $x \in |X|$, and set
$y = f(x) \in |Y|$.

\medskip\noindent
The case $\mathcal{P} = (\beta)$. Condition $(\beta)$ for $Y$ means
there exists a monomorphism $\Spec(k) \to Y$ representing $y$.
The fibre product $X_y = \Spec(k) \times_Y X$ is a scheme, and
$x$ corresponds to a point of $X_y$, i.e., to a monomorphism
$\Spec(k') \to X_y$. As $X_y \to X$ is a monomorphism also we see
that $x$ is represented by the monomorphism $\Spec(k') \to X_y \to X$.
In other words $(\beta)$ holds for $X$.

\medskip\noindent
The case $\mathcal{P} = (\gamma)$. Since $(\gamma) \Rightarrow (\beta)$
we have seen in the preceding paragraph that $y$ and $x$ can be represented
by monomorphisms as in the following diagram
$$
\xymatrix{
\Spec(k') \ar[r]_-x \ar[d] & X \ar[d] \\
\Spec(k) \ar[r]^-y & Y
}
$$
Also, by definition of property $(\gamma)$ via
Lemma \ref{lemma-UR-finite-above-x} (2)
there exist schemes
$V_i$ and \'etale morphisms $V_i \to Y$ such that $\coprod V_i \to Y$
is surjective and for each $i$, setting $R_i = V_i \times_Y V_i$
the fibres of both
$$
|V_i| \longrightarrow |Y|
\quad\text{and}\quad
|R_i| \longrightarrow |Y|
$$
over $y$ are finite. This means that the schemes
$(V_i)_y$ and $(R_i)_y$ are finite schemes over $y = \Spec(k)$.
As $X \to Y$ is representable, the fibre products $U_i = V_i \times_Y X$
are schemes. The morphisms $U_i \to X$ are \'etale, and
$\coprod U_i \to X$ is surjective. Finally, for each $i$ we have
$$
(U_i)_x =
(V_i \times_Y X)_x =
(V_i)_y \times_{\Spec(k)} \Spec(k')
$$
and
$$
(U_i \times_X U_i)_x =
\left((V_i \times_Y X) \times_X (V_i \times_Y X)\right)_x =
(R_i)_y \times_{\Spec(k)} \Spec(k')
$$
hence these are finite over $k'$ as base changes of the finite
schemes $(V_i)_y$ and $(R_i)_y$. This implies that $(\gamma)$ holds for $X$,
again via the second condition of
Lemma \ref{lemma-UR-finite-above-x}.

\medskip\noindent
The case $\mathcal{P} = (\delta)$. Let $V \to Y$ be an \'etale morphism with
$V$ an affine scheme. Since $Y$ has property $(\delta)$ this morphism has
universally bounded fibres. By
Lemma \ref{lemma-base-change-universally-bounded}
the base change $V \times_Y X \to X$ also has universally bounded fibres.
Hence the first part of
Lemma \ref{lemma-U-universally-bounded}
applies and we see that $Y$ also has property $(\delta)$.

\medskip\noindent
The case $\mathcal{P} = (\epsilon)$. We will repeatedly use
Spaces, Lemma
\ref{spaces-lemma-base-change-representable-transformations-property}.
Let $V_i \to Y$ be as in
Lemma \ref{lemma-characterize-very-reasonable} (2).
Set $U_i = X \times_Y V_i$. The morphisms $U_i \to X$ are \'etale,
and $\coprod U_i \to X$ is surjective. Because
$U_i \times_X U_i = X \times_Y (V_i \times_Y V_i)$ we see
that the projections $U_i \times_Y U_i \to U_i$ are
base changes of the projections $V_i \times_Y V_i \to V_i$, and so
quasi-compact as well. Hence $X$ satisfies
Lemma \ref{lemma-characterize-very-reasonable} (2).

\medskip\noindent
The case $\mathcal{P} = (\theta)$. In this case the result is
Categories, Lemma \ref{categories-lemma-representable-over-representable}.
\end{proof}

\end{document}
